% page 53 of the facsimile (image p-052 of the scan)
% block 165.1 x 242.0 mm ; left margin 36.9 mm
\pgc{15.240mm}{0.000mm}{
\l{\cel{51}45}
\l{}
\l{\sou{asperulo}. (\sou{bot.}) Genero de rubiacei de la regioni ye temperaturo}
\l{\cel{3}meza. - L.\cel{1}\sou{asperula}.\cel{2}- EFIS.}
\l{}
\l{\sou{aspiko}. (\sou{zool.}) Varietato di vipero, serpento venenoza di Egiptia}
\l{\cel{3}e di Libia, quan karakterizas influro remarkinda di lua kolo.}
\l{\cel{3}- L.\cel{1}\sou{naja haje}. - EFIRS.}
\l{}
\l{\sou{aspikjeleo}.- Karno-disho o fisho-disho kun jeleo. - DEFIS.}
\l{}
\l{\sou{aspirar}. (\sou{trans, per, ad}). I. Atraktar per inspiro di aero. -}
\l{\cel{3}II. Direktar sua deziro ad ulu od ulo. - DEFIRS.}
\l{}
\l{\sou{aspiratoro}. - Nomo di aparati diversa qui havas kom rolo aspirar}
\l{\cel{3}la fluidi (aero, gaso, vapori, liquidi, e c.) pura o qui}
\l{\cel{3}kontenas korpi solida suspende, la polvi (segala, limala,}
\l{\cel{3}e c.) - DEFIRS.}
\l{}
\l{\sou{aspirino}. Acido acetil-salicila : CH3 COOC6H4COOH. Korpo blanka,}
\l{\cel{3}solida, fuzo-punto : 133o Celsius. Uzata medicine kom kontre-}
\l{\cel{3}febro, kontre-doloro. - DEFIRS.}
\l{}
\l{\sou{astenopa}. (\sou{patol.}) Dicesas pri persono di qua la vidiveso ne}
\l{\cel{3}esas konstanta. (Komence, la malado vidas bone ; balde, lu}
\l{\cel{3}vidas quaze tra nebulo ; fine, lu ne plus povas dicernar la}
\l{\cel{3}objekti). - DEFIS.}
\l{}
\l{\sou{astatika}. (\sou{magnet.}) Qua permanas movebla kande deprenata de sua}
\l{\cel{3}cirkonstanci naturala. - DEFIRS.}
\l{}
\l{\sou{astero}. (\sou{bot.})\cel{2}Planto kompozaja di qua la flori radioza simil-}
\l{\cel{3}esas steli. - L.\cel{1}\sou{aster}. - DEFIRS.}
\l{}
\l{\sou{asterio}.- Zoofito radiera. - L.\cel{1}\sou{asterias}. - DEFIRS.}
\l{}
\l{\sou{asterismo}. Fenomeno optikala : reflekto radiacanta di lumo}
\l{\cel{3}sur ula lapidi. - DEFIRS.}
\l{}
\l{\sou{asteroido}. (\sou{astron.})\cel{2}Mikra korpo qua cirkulas en cielo.- DEFIRS.}
\l{}
\l{\sou{astigmata}. (\sou{patol.})\cel{2}Di qua la okulo havas kom karakterizivo}
\l{\cel{3}kurveso-neegalaji en la medii refraktiva. - DEF.}
\l{}
\l{\sou{astmo}. (\sou{patol.})Sufoko intermitanta, atribuata a nevroso di}
\l{\cel{3}la respir-organi. - DEFIRS.}
\l{}
\l{\sou{astonar}. (\sou{trans.}) Shokar la psiko per ulo extraordinara. - EF.}
\l{}
\l{\sou{astragalo}. I. (\sou{anat.}) Osto di la tarso, di formo kuboida, in-}
\l{\cel{3}kluzita inter la maleoli di la tibio e di la peroneo. -}
\l{\cel{3}II. (\sou{arkitekt.})\cel{2}Muluro pozita sur kolono, inter la fusto}
\l{\cel{3}e la kapitelo. - DEFIRS.}
\l{}
\l{\sou{astrakano}. Felo di qua la pili esas lokloza, devenanta de}
\l{\cel{3}bovido nasko-mortinta. - DEFIRS.}
}
